\chapter{Kernreaktionen, Spaltung und Fusion}

\section{Ziel dieses Kapitels}

Dieses Kapitel behandelt Kernreaktionen.
Eine Kernreaktion ist ein Prozess, bei dem ein Projektil mit einem Kern wechselwirkt
und dadurch neue Kerne, Teilchen oder angeregte Zustaende entstehen.

Im Zentrum stehen:
\begin{itemize}
    \item Reaktionsgleichungen,
    \item Erhaltungssaetze,
    \item Q-Werte,
    \item Schwellenenergien,
    \item Wirkungsquerschnitte,
    \item Kernspaltung,
    \item Kettenreaktionen,
    \item Kernfusion,
    \item Energieerzeugung in Sternen,
    \item Bragg-Peak und Energieverlust in Materie.
\end{itemize}

\begin{notebox}
Dieses Kapitel behandelt die Stichwoerter:
\begin{itemize}
    \item Kernreaktionen,
    \item Q-Wert,
    \item Reaktionsenergie,
    \item Schwellenenergie,
    \item Wirkungsquerschnitt,
    \item Compound-Kern,
    \item Spaltung,
    \item induzierte Spaltung,
    \item spontane Spaltung,
    \item Kettenreaktion,
    \item kritische Masse,
    \item Kernfusion,
    \item Coulomb-Barriere,
    \item Tunneleffekt,
    \item Proton-Proton-Kette,
    \item CNO-Zyklus,
    \item Bragg-Peak.
\end{itemize}
\end{notebox}

\section{Was ist eine Kernreaktion?}

\begin{definitionbox}
Eine \emph{Kernreaktion} ist eine Wechselwirkung, bei der ein Atomkern durch ein
einfallendes Teilchen oder einen anderen Kern veraendert wird.
\end{definitionbox}

\begin{conceptbox}
Bei einer Kernreaktion koennen verschiedene Dinge passieren:
\begin{itemize}
    \item ein Kern wird angeregt,
    \item ein Teilchen wird elastisch gestreut,
    \item ein Teilchen wird inelastisch gestreut,
    \item ein Teilchen wird eingefangen,
    \item ein Kern wird gespalten,
    \item zwei Kerne verschmelzen,
    \item neue Teilchen werden erzeugt.
\end{itemize}
\end{conceptbox}

\begin{formulabox}
Eine typische Kernreaktion schreibt man als
\[
    a + A \rightarrow b + B.
\]
Dabei ist:
\begin{itemize}
    \item $a$ das einfallende Projektil,
    \item $A$ der Targetkern,
    \item $b$ ein auslaufendes Teilchen,
    \item $B$ der Restkern beziehungsweise Produktkern.
\end{itemize}
\end{formulabox}

\begin{notebox}
Man verwendet auch die Kurzschreibweise:
\[
    A(a,b)B.
\]
Diese bedeutet dasselbe wie:
\[
    a + A \rightarrow b + B.
\]
\end{notebox}

\begin{examplebox}
Die Reaktion
\[
    {}^{14}\mathrm{N}(\alpha,p){}^{17}\mathrm{O}
\]
bedeutet ausgeschrieben:
\[
    \alpha + {}^{14}\mathrm{N}
    \rightarrow
    p + {}^{17}\mathrm{O}.
\]
Ein Alpha-Teilchen trifft auf Stickstoff-14.
Am Ende entstehen ein Proton und Sauerstoff-17.
\end{examplebox}

\section{Erhaltungssaetze bei Kernreaktionen}

Kernreaktionen muessen grundlegende Erhaltungssaetze erfuellen.

\begin{notebox}
Wichtige Erhaltungsgroessen sind:
\begin{itemize}
    \item elektrische Ladung,
    \item Nukleonenzahl beziehungsweise Baryonenzahl,
    \item Energie,
    \item Impuls,
    \item Drehimpuls,
    \item Paritaet, falls die Wechselwirkung sie erhaelt.
\end{itemize}
\end{notebox}

\begin{conceptbox}
In einfachen kernphysikalischen Reaktionsgleichungen prueft man zuerst:
\[
    A_{\mathrm{links}} = A_{\mathrm{rechts}},
\]
\[
    Z_{\mathrm{links}} = Z_{\mathrm{rechts}}.
\]

Das bedeutet:
Massenzahl und Ladungszahl muessen erhalten sein.
\end{conceptbox}

\begin{examplebox}
Betrachte:
\[
    {}^{2}_{1}\mathrm{H}
    +
    {}^{3}_{1}\mathrm{H}
    \rightarrow
    {}^{4}_{2}\mathrm{He}
    +
    n.
\]

Massenzahlen:
\[
    2+3 = 4+1.
\]

Protonenzahlen:
\[
    1+1 = 2+0.
\]

Die Reaktion erhaelt also Massenzahl und elektrische Ladung.
\end{examplebox}

\begin{examtipbox}
Bei jeder Kernreaktion zuerst kontrollieren:
\begin{itemize}
    \item Ist die Summe der Massenzahlen $A$ erhalten?
    \item Ist die Summe der Protonenzahlen $Z$ erhalten?
    \item Ist der Q-Wert positiv oder negativ?
\end{itemize}
\end{examtipbox}

\section{Q-Wert einer Kernreaktion}

Der Q-Wert beschreibt, ob bei einer Kernreaktion Energie frei wird oder Energie
zugefuehrt werden muss.

\begin{definitionbox}
Der \emph{Q-Wert} einer Kernreaktion ist die Differenz der Ruheenergien vor und
nach der Reaktion:
\[
    Q =
    \left(
        m_{\mathrm{Anfang}} - m_{\mathrm{Ende}}
    \right)c^2.
\]
\end{definitionbox}

\begin{formulabox}
Fuer
\[
    a + A \rightarrow b + B
\]
gilt:
\[
    Q =
    \left(
        m_a + m_A - m_b - m_B
    \right)c^2.
\]
\end{formulabox}

\begin{conceptbox}
Wenn
\[
    Q>0
\]
ist, wird Energie frei.
Die Reaktion ist exoenergetisch.

Wenn
\[
    Q<0
\]
ist, muss Energie zugefuehrt werden.
Die Reaktion ist endoenergetisch.
\end{conceptbox}

\begin{examtipbox}
Ein positiver Q-Wert bedeutet nicht automatisch, dass die Reaktion bei beliebig
kleiner Projektilenergie schnell ablaeuft.

Auch bei positivem Q-Wert koennen Coulomb-Barrieren, Drehimpulsbarrieren oder kleine
Wirkungsquerschnitte die Reaktion stark behindern.
\end{examtipbox}

\section{Schwellenenergie}

Wenn eine Reaktion einen negativen Q-Wert hat, muss mindestens eine bestimmte
kinetische Energie zugefuehrt werden.

\begin{definitionbox}
Die \emph{Schwellenenergie} ist die minimale Projektilenergie, die noetig ist,
damit eine endoenergetische Reaktion ablaufen kann.
\end{definitionbox}

\begin{conceptbox}
Wenn
\[
    Q<0
\]
ist, muss mindestens die Energie
\[
    |Q|
\]
in die Reaktion eingebracht werden.

Im Laborsystem ist die notwendige Projektilenergie meist groesser als $|Q|$,
weil Impuls erhalten bleiben muss und ein Teil der Energie in Bewegung des
Schwerpunktsystems steckt.
\end{conceptbox}

\begin{formulabox}
Fuer eine Reaktion
\[
    a + A \rightarrow b + B
\]
mit ruhendem Target $A$ ist die Schwellenenergie im nichtrelativistischen Grenzfall
naeherungsweise:
\[
    E_{\mathrm{thr}}
    \approx
    -Q
    \left(
        1+\frac{m_a}{m_A}
    \right)
\]
fuer $Q<0$.
\end{formulabox}

\begin{notebox}
Die genaue Schwellenenergie folgt aus relativistischer Energie- und Impulserhaltung.
Die obige Formel ist eine nuetzliche Naeherung fuer viele kernphysikalische Aufgaben.
\end{notebox}

\section{Wirkungsquerschnitt von Kernreaktionen}

Wie wahrscheinlich eine Kernreaktion ist, wird durch den Wirkungsquerschnitt
beschrieben.

\begin{definitionbox}
Der \emph{Wirkungsquerschnitt} $\sigma$ ist ein Mass fuer die Wahrscheinlichkeit,
dass eine bestimmte Reaktion stattfindet.
\end{definitionbox}

\begin{conceptbox}
Der Wirkungsquerschnitt haengt im Allgemeinen ab von:
\begin{itemize}
    \item der Energie des Projektils,
    \item der Art des Projektils,
    \item dem Targetkern,
    \item der Wechselwirkung,
    \item moeglichen Resonanzen,
    \item Drehimpuls- und Paritaetsauswahlregeln.
\end{itemize}
\end{conceptbox}

\begin{formulabox}
Die Reaktionsrate ist grob proportional zu:
\[
    R \propto \Phi \, N_{\mathrm{Target}} \, \sigma.
\]
Dabei ist:
\begin{itemize}
    \item $\Phi$ der Teilchenfluss,
    \item $N_{\mathrm{Target}}$ die Zahl der Targetkerne,
    \item $\sigma$ der Wirkungsquerschnitt.
\end{itemize}
\end{formulabox}

\begin{examtipbox}
Ein grosser Q-Wert bedeutet nicht automatisch einen grossen Wirkungsquerschnitt.

Der Q-Wert sagt etwas ueber die Energiebilanz.
Der Wirkungsquerschnitt sagt etwas ueber die Wahrscheinlichkeit des Prozesses.
\end{examtipbox}

\section{Compound-Kern}

Viele Kernreaktionen laufen ueber einen Zwischenzustand ab.
Diesen Zwischenzustand nennt man Compound-Kern.

\begin{definitionbox}
Ein \emph{Compound-Kern} ist ein angeregter Zwischenkern, der entsteht, wenn ein
Projektil vom Targetkern eingefangen wird.
\end{definitionbox}

\begin{formulabox}
Schematisch:
\[
    a + A \rightarrow C^* \rightarrow b + B.
\]
Dabei ist $C^*$ der angeregte Compound-Kern.
\end{formulabox}

\begin{conceptbox}
Die Energie des einfallenden Projektils verteilt sich im Compound-Kern auf viele
Nukleonen.

Der Compound-Kern verliert anschliessend Energie durch:
\begin{itemize}
    \item Aussenden eines Teilchens,
    \item Gamma-Emission,
    \item Spaltung.
\end{itemize}
\end{conceptbox}

\section{Kernspaltung}

Bei der Kernspaltung zerfaellt ein schwerer Kern in zwei mittelschwere Fragmente.

\begin{definitionbox}
\emph{Kernspaltung} ist die Aufteilung eines schweren Atomkerns in zwei leichtere
Kerne, meist begleitet von der Emission mehrerer Neutronen und Gamma-Strahlung.
\end{definitionbox}

\begin{conceptbox}
Kernspaltung ist energetisch moeglich, weil sehr schwere Kerne eine kleinere
Bindungsenergie pro Nukleon besitzen als mittelschwere Kerne.

Wenn ein schwerer Kern in mittelschwere Fragmente zerfaellt, steigt die
Bindungsenergie pro Nukleon.
Die Differenz wird als Energie frei.
\end{conceptbox}

\begin{formulabox}
Schematisch:
\[
    {}^{A}_{Z}X
    \rightarrow
    {}^{A_1}_{Z_1}Y
    +
    {}^{A_2}_{Z_2}W
    +
    k n
    +
    Q.
\]
Dabei gilt:
\[
    A=A_1+A_2+k,
\]
\[
    Z=Z_1+Z_2.
\]
\end{formulabox}

\begin{examtipbox}
Bei Spaltungsreaktionen immer beachten:
\begin{itemize}
    \item Massenzahl bleibt insgesamt erhalten.
    \item Ladungszahl bleibt erhalten.
    \item Mehrere Neutronen koennen frei werden.
    \item Der Q-Wert ist bei schweren Kernen typischerweise positiv.
\end{itemize}
\end{examtipbox}

\section{Induzierte Spaltung}

Viele schwere Kerne koennen durch Neutroneneinfang zur Spaltung angeregt werden.

\begin{definitionbox}
Bei der \emph{induzierten Spaltung} wird ein schwerer Kern durch ein einfallendes
Teilchen, meist ein Neutron, zur Spaltung gebracht.
\end{definitionbox}

\begin{examplebox}
Ein klassisches Beispiel ist:
\[
    n + {}^{235}_{92}\mathrm{U}
    \rightarrow
    {}^{236}_{92}\mathrm{U}^{*}
    \rightarrow
    \text{Spaltfragmente}
    + \text{Neutronen}
    + Q.
\]
Das Neutron wird eingefangen.
Es entsteht ein angeregter Compound-Kern ${}^{236}\mathrm{U}^{*}$, der spalten kann.
\end{examplebox}

\begin{conceptbox}
Neutronen sind als Projektile fuer Spaltung besonders effektiv, weil sie keine
elektrische Ladung besitzen.
Sie muessen daher keine Coulomb-Barriere ueberwinden.
\end{conceptbox}

\section{Spontane Spaltung}

\begin{definitionbox}
Bei der \emph{spontanen Spaltung} spaltet ein schwerer Kern ohne ein einfallendes
Projektil.
\end{definitionbox}

\begin{conceptbox}
Spontane Spaltung ist ein quantenmechanischer Prozess.
Der Kern kann durch eine Spaltbarriere tunneln.

Sie wird besonders relevant bei sehr schweren Kernen, bei denen die Coulomb-Abstossung
gross ist.
\end{conceptbox}

\begin{notebox}
In der Nuklidkarte werden bei sehr schweren Kernen Zerfallstypen wie Alpha-Zerfall
und spontane Spaltung markiert.
Die Konkurrenz dieser Zerfallsarten ist wichtig fuer superschwere Elemente.
\end{notebox}

\section{Kettenreaktion}

Bei einer Spaltung werden Neutronen frei.
Diese Neutronen koennen weitere Spaltungen ausloesen.

\begin{definitionbox}
Eine \emph{Kettenreaktion} ist eine Folge von Spaltungen, bei der die bei einer
Spaltung frei werdenden Neutronen weitere Spaltungen ausloesen.
\end{definitionbox}

\begin{conceptbox}
Wenn im Mittel weniger als ein Neutron pro Spaltung eine weitere Spaltung ausloest,
stirbt die Kettenreaktion ab.

Wenn im Mittel genau ein Neutron pro Spaltung eine weitere Spaltung ausloest,
ist die Reaktion kritisch.

Wenn im Mittel mehr als ein Neutron pro Spaltung eine weitere Spaltung ausloest,
waechst die Reaktion an.
\end{conceptbox}

\begin{definitionbox}
Der \emph{Multiplikationsfaktor} $k$ beschreibt, wie viele Neutronen einer Generation
im Mittel zur naechsten Generation beitragen.
\end{definitionbox}

\begin{formulabox}
\[
    k<1
    \quad\Rightarrow\quad
    \text{unterkritisch}.
\]
\[
    k=1
    \quad\Rightarrow\quad
    \text{kritisch}.
\]
\[
    k>1
    \quad\Rightarrow\quad
    \text{ueberkritisch}.
\]
\end{formulabox}

\section{Kritische Masse und Moderation}

\begin{definitionbox}
Die \emph{kritische Masse} ist die minimale Masse eines spaltbaren Materials,
bei der eine selbsterhaltende Kettenreaktion moeglich ist.
\end{definitionbox}

\begin{conceptbox}
Ob eine Kettenreaktion erhalten bleibt, haengt nicht nur von der Masse ab.
Wichtig sind auch:
\begin{itemize}
    \item Form des Materials,
    \item Dichte,
    \item Neutronenverluste nach aussen,
    \item Neutronenabsorption,
    \item Energie der Neutronen,
    \item Moderation.
\end{itemize}
\end{conceptbox}

\begin{definitionbox}
Ein \emph{Moderator} ist ein Material, das schnelle Neutronen abbremst.
\end{definitionbox}

\begin{conceptbox}
Bei manchen spaltbaren Materialien ist die Spaltwahrscheinlichkeit fuer langsame
Neutronen groesser als fuer schnelle Neutronen.

Ein Moderator kann daher die Kettenreaktion beguenstigen, indem er Neutronen abbremst.
\end{conceptbox}

\section{Verzoegerte Neutronen}

\begin{definitionbox}
\emph{Verzoegerte Neutronen} sind Neutronen, die nicht direkt bei der Spaltung,
sondern erst beim radioaktiven Zerfall bestimmter Spaltprodukte entstehen.
\end{definitionbox}

\begin{conceptbox}
Verzoegerte Neutronen sind fuer die Steuerung von Kernreaktoren sehr wichtig.

Ohne verzoegerte Neutronen wuerde sich die Neutronenzahl extrem schnell aendern.
Mit verzoegerten Neutronen wird die Reaktordynamik langsamer und kontrollierbarer.
\end{conceptbox}

\begin{examtipbox}
Pruefungsrelevanter Merksatz:
Verzoegerte Neutronen sind entscheidend fuer die technische Kontrollierbarkeit eines
Kernreaktors.
\end{examtipbox}

\section{Energie bei der Spaltung}

\begin{conceptbox}
Die Energie aus der Spaltung stammt aus dem Massendefekt.
Die Spaltprodukte sind insgesamt staerker gebunden als der urspruengliche schwere
Kern.

Die freiwerdende Energie erscheint hauptsaechlich als:
\begin{itemize}
    \item kinetische Energie der Spaltfragmente,
    \item kinetische Energie der Neutronen,
    \item Gamma-Strahlung,
    \item Beta-Zerfallsenergie der Spaltprodukte,
    \item Neutrinoenergie.
\end{itemize}
\end{conceptbox}

\begin{notebox}
Ein typischer Richtwert fuer die freiwerdende Energie pro Spaltung eines schweren
Kerns ist etwa
\[
    \SI{200}{MeV}.
\]
Der genaue Wert haengt vom spaltenden Kern und den Spaltfragmenten ab.
\end{notebox}

\section{Kernfusion}

Bei der Kernfusion verschmelzen leichte Kerne zu einem schwereren Kern.

\begin{definitionbox}
\emph{Kernfusion} ist die Verschmelzung leichter Atomkerne zu einem schwereren Kern.
\end{definitionbox}

\begin{conceptbox}
Fusion leichter Kerne setzt Energie frei, weil die Bindungsenergie pro Nukleon
bei leichten Kernen mit wachsender Massenzahl stark zunimmt.

Die Produkte sind staerker gebunden als die Ausgangskerne.
Die Massendifferenz wird als Energie frei.
\end{conceptbox}

\begin{examplebox}
Eine wichtige Fusionsreaktion ist:
\[
    {}^{2}_{1}\mathrm{H}
    +
    {}^{3}_{1}\mathrm{H}
    \rightarrow
    {}^{4}_{2}\mathrm{He}
    +
    n
    +
    Q.
\]
Diese Reaktion ist die Deuterium-Tritium-Fusion.
\end{examplebox}

\section{Coulomb-Barriere bei der Fusion}

Kerne sind positiv geladen.
Damit zwei Kerne fusionieren koennen, muessen sie sich sehr nahe kommen.
Dabei muessen sie die elektrische Abstossung ueberwinden.

\begin{definitionbox}
Die \emph{Coulomb-Barriere} ist die Energiebarriere, die durch die elektrische
Abstossung zweier positiv geladener Kerne entsteht.
\end{definitionbox}

\begin{formulabox}
Fuer zwei Kerne mit Ladungen $Z_1e$ und $Z_2e$ ist das Coulomb-Potential
naeherungsweise:
\[
    V_C(r)
    =
    \frac{1}{4\pi\varepsilon_0}
    \frac{Z_1Z_2e^2}{r}.
\]
\end{formulabox}

\begin{conceptbox}
Klassisch muessten die Kerne genug kinetische Energie besitzen, um die Coulomb-Barriere
zu ueberwinden.

Quantenmechanisch koennen sie aber mit endlicher Wahrscheinlichkeit durch die Barriere
tunneln.

Dieser Tunneleffekt ist entscheidend fuer Kernfusion in Sternen.
\end{conceptbox}

\begin{examtipbox}
Fusion in Sternen waere klassisch bei den dortigen Temperaturen zu selten.
Der quantenmechanische Tunneleffekt macht Fusion trotzdem moeglich.
\end{examtipbox}

\section{Proton-Proton-Kette}

Die Proton-Proton-Kette ist ein wichtiger Fusionsprozess in der Sonne.

\begin{definitionbox}
Die \emph{Proton-Proton-Kette} ist eine Folge von Fusionsreaktionen, bei der aus
Wasserstoff letztlich Helium entsteht.
\end{definitionbox}

\begin{conceptbox}
Die Netto-Reaktion der Wasserstofffusion in der Sonne ist:
\[
    4\,{}^{1}\mathrm{H}
    \rightarrow
    {}^{4}\mathrm{He}
    + 2e^+
    + 2\nu_e
    + Q.
\]
Beruecksichtigt man die Annihilation der Positronen mit Elektronen, wird pro
Netto-Reaktion eine Energie von etwa
\[
    Q \approx \SI{26.7}{MeV}
\]
frei.
\end{conceptbox}

\begin{notebox}
Die Vorlesungsabbildung zur Proton-Proton-Kette zeigt verschiedene Aeste.
Ein wichtiger erster Schritt ist:
\[
    {}^{1}\mathrm{H}
    +
    {}^{1}\mathrm{H}
    \rightarrow
    {}^{2}\mathrm{H}
    +
    e^+
    +
    \nu_e.
\]
Danach entsteht ueber weitere Reaktionen Helium.
\end{notebox}

\begin{conceptbox}
Der erste Schritt der pp-Kette ist schwach wechselwirkend, weil ein Proton in ein
Neutron umgewandelt werden muss.

Deshalb ist dieser Schritt langsam.
Das ist wichtig fuer die lange Lebensdauer der Sonne.
\end{conceptbox}

\section{CNO-Zyklus}

Neben der Proton-Proton-Kette gibt es den CNO-Zyklus.
Er ist besonders in massereicheren Sternen wichtig.

\begin{definitionbox}
Der \emph{CNO-Zyklus} ist eine Fusionsreaktionskette, bei der Kohlenstoff, Stickstoff
und Sauerstoff als Katalysatoren fuer die Umwandlung von Wasserstoff in Helium wirken.
\end{definitionbox}

\begin{conceptbox}
Die Netto-Reaktion des CNO-Zyklus ist ebenfalls:
\[
    4\,{}^{1}\mathrm{H}
    \rightarrow
    {}^{4}\mathrm{He}
    + 2e^+
    + 2\nu_e
    + Q.
\]

Kohlenstoff, Stickstoff und Sauerstoff werden dabei im Zyklus verbraucht und wieder
erzeugt.
Sie wirken daher als Katalysatoren.
\end{conceptbox}

\begin{notebox}
Die Vorlesungsabbildung zum CNO-Zyklus zeigt Reaktionen wie:
\[
    {}^{12}\mathrm{C}
    +
    {}^{1}\mathrm{H}
    \rightarrow
    {}^{13}\mathrm{N}
    +
    \gamma.
\]
Danach folgen Beta-Zerfaelle und weitere Protoneneinfangreaktionen, bis am Ende
wieder ${}^{12}\mathrm{C}$ entsteht.
\end{notebox}

\begin{examtipbox}
Proton-Proton-Kette und CNO-Zyklus haben dieselbe Netto-Reaktion:
Wasserstoff wird zu Helium fusioniert.

Der Unterschied ist:
\begin{itemize}
    \item pp-Kette: dominiert in sonnenaehnlichen Sternen.
    \item CNO-Zyklus: wichtiger bei hoeheren Temperaturen und massereicheren Sternen.
\end{itemize}
\end{examtipbox}

\section{Fusion und Temperatur}

\begin{conceptbox}
Fusion benoetigt hohe Temperaturen, weil die Kerne genug kinetische Energie brauchen,
um nahe genug zusammenzukommen.

Die Temperatur bestimmt die Geschwindigkeitsverteilung der Kerne.
Nur ein Teil der Kerne besitzt genug Energie, um mit relevanter Wahrscheinlichkeit
zu fusionieren.
\end{conceptbox}

\begin{definitionbox}
Das \emph{Gamow-Fenster} ist der Energiebereich, in dem Fusionsreaktionen in Sternen
hauptsaechlich stattfinden.

Es ergibt sich aus dem Zusammenspiel von thermischer Energieverteilung und
Tunnelwahrscheinlichkeit durch die Coulomb-Barriere.
\end{definitionbox}

\begin{conceptbox}
Bei niedriger Energie gibt es viele Teilchen, aber die Tunnelwahrscheinlichkeit ist
sehr klein.

Bei hoher Energie ist die Tunnelwahrscheinlichkeit groesser, aber es gibt nur wenige
Teilchen mit dieser Energie.

Das Maximum des Produkts liegt im Gamow-Fenster.
\end{conceptbox}

\section{Bragg-Peak und Energieverlust in Materie}

Geladene Teilchen verlieren beim Durchgang durch Materie Energie.
Dieser Energieverlust ist fuer Kernreaktionen, Strahlenschutz und medizinische
Anwendungen wichtig.

\begin{definitionbox}
Der \emph{Bragg-Peak} ist das Maximum der Energieabgabe eines geladenen Teilchens
kurz vor dem Ende seiner Reichweite in Materie.
\end{definitionbox}

\begin{conceptbox}
Ein geladenes Teilchen verliert Energie vor allem durch Ionisation und Anregung
von Atomen.

Wenn das Teilchen langsamer wird, steigt der Energieverlust pro Weglaenge zunaechst an.
Kurz vor dem Stillstand tritt ein Maximum auf: der Bragg-Peak.
\end{conceptbox}

\begin{notebox}
Die Vorlesungsabbildung zur Ionisation von Alpha-Teilchen zeigt genau dieses Verhalten:
Die spezifische Ionisation steigt entlang der Bahn und erreicht kurz vor dem Ende
der Reichweite ein Maximum.
\end{notebox}

\begin{examtipbox}
Der Bragg-Peak ist wichtig, weil geladene Teilchen ihre Energie sehr lokal am Ende
ihrer Reichweite abgeben koennen.

Das ist zum Beispiel fuer Teilchentherapie relevant.
\end{examtipbox}

\section{Erzeugung schwerer und superschwerer Kerne}

Sehr schwere Kerne koennen durch Kernreaktionen erzeugt werden.
Dabei werden schwere Ionen auf Targetkerne geschossen.

\begin{conceptbox}
Um ein neues schweres Element zu erzeugen, muss ein Projektilkern mit einem Targetkern
verschmelzen.

Dabei entsteht ein hochangeregter Compound-Kern.
Dieser kann entweder durch Abdampfen von Neutronen abkuehlen oder wieder zerfallen.
\end{conceptbox}

\begin{formulabox}
Schematisch:
\[
    {}^{A_1}_{Z_1}X
    +
    {}^{A_2}_{Z_2}Y
    \rightarrow
    {}^{A_1+A_2}_{Z_1+Z_2}C^{*}
    \rightarrow
    \text{Produktkern}
    +
    x n.
\]
\end{formulabox}

\begin{conceptbox}
Die Erzeugung superschwerer Elemente ist schwierig, weil:
\begin{itemize}
    \item die Coulomb-Barriere sehr hoch ist,
    \item der Compound-Kern stark angeregt ist,
    \item Spaltung mit der Abkuehlung durch Neutronenemission konkurriert,
    \item die Wirkungsquerschnitte sehr klein sind.
\end{itemize}
\end{conceptbox}

\section{Vergleich: Spaltung und Fusion}

\begin{center}
\begin{tabular}{|p{4.2cm}|p{5.1cm}|p{5.1cm}|}
\hline
Eigenschaft & Spaltung & Fusion \\
\hline
Ausgangskerne & schwere Kerne & leichte Kerne \\
\hline
Produkte & mittelschwere Kerne und Neutronen & schwererer Kern, eventuell Teilchen \\
\hline
Energiegrund & Produkte staerker gebunden & Produkt staerker gebunden \\
\hline
Hindernis & Spaltbarriere, Neutronenbilanz & Coulomb-Barriere \\
\hline
Typisches Beispiel & Spaltung von Uran oder Plutonium & Wasserstofffusion zu Helium \\
\hline
\end{tabular}
\end{center}

\begin{examtipbox}
Spaltung und Fusion setzen beide Energie frei, wenn die Endprodukte staerker gebunden
sind als die Anfangskerne.

Der Unterschied ist:
\begin{itemize}
    \item Fusion geht bei leichten Kernen in Richtung hoeherer Bindungsenergie pro Nukleon.
    \item Spaltung geht bei sehr schweren Kernen in Richtung hoeherer Bindungsenergie pro Nukleon.
\end{itemize}
\end{examtipbox}

\section{Mini-Zusammenfassung}

\begin{notebox}
Die wichtigsten Punkte dieses Kapitels sind:
\begin{itemize}
    \item Kernreaktionen muessen Ladung, Nukleonenzahl, Energie und Impuls erhalten.
    \item Der Q-Wert ist
    \[
        Q =
        \left(
            m_{\mathrm{Anfang}}
            -
            m_{\mathrm{Ende}}
        \right)c^2.
    \]
    \item Bei $Q>0$ wird Energie frei.
    \item Bei $Q<0$ braucht man eine Schwellenenergie.
    \item Der Wirkungsquerschnitt beschreibt die Wahrscheinlichkeit einer Reaktion.
    \item Ein Compound-Kern ist ein angeregter Zwischenkern.
    \item Spaltung schwerer Kerne setzt Energie frei, weil die Fragmente staerker gebunden sind.
    \item Eine Kettenreaktion entsteht durch bei Spaltung frei werdende Neutronen.
    \item Fusion leichter Kerne setzt Energie frei, weil das Produkt staerker gebunden ist.
    \item Fusion muss die Coulomb-Barriere ueberwinden oder durchtunneln.
    \item Die Proton-Proton-Kette und der CNO-Zyklus wandeln netto Wasserstoff in Helium um.
    \item Der Bragg-Peak beschreibt maximale Energieabgabe kurz vor dem Stillstand geladener Teilchen.
\end{itemize}
\end{notebox}

\section{Pruefungscheck}

\begin{examtipbox}
Nach diesem Kapitel solltest du folgende Fragen beantworten koennen:
\begin{enumerate}
    \item Was ist eine Kernreaktion?
    \item Wie liest man die Schreibweise $A(a,b)B$?
    \item Welche Erhaltungssaetze gelten bei Kernreaktionen?
    \item Was ist der Q-Wert?
    \item Was bedeutet $Q>0$?
    \item Was bedeutet $Q<0$?
    \item Was ist eine Schwellenenergie?
    \item Was beschreibt der Wirkungsquerschnitt?
    \item Was ist ein Compound-Kern?
    \item Was ist Kernspaltung?
    \item Warum setzt Spaltung schwerer Kerne Energie frei?
    \item Was ist induzierte Spaltung?
    \item Warum sind Neutronen gute Projektile fuer Spaltung?
    \item Was ist eine Kettenreaktion?
    \item Was bedeuten unterkritisch, kritisch und ueberkritisch?
    \item Warum sind verzoegerte Neutronen wichtig?
    \item Was ist Kernfusion?
    \item Warum gibt es bei Fusion eine Coulomb-Barriere?
    \item Warum ist der Tunneleffekt fuer Sternfusion wichtig?
    \item Was ist die Netto-Reaktion der Proton-Proton-Kette?
    \item Was ist der CNO-Zyklus?
    \item Was ist der Bragg-Peak?
    \item Warum ist die Erzeugung superschwerer Elemente schwierig?
\end{enumerate}
\end{examtipbox}